\documentclass{situationist-preprint}

\title{The Responsibility of Groups}
\msid{ACT-10-b}
\msfamily{ACT-10}
\msdiscipline{Action, Agency and Free Will}
\mstype{Companion manuscript}
\mskeywords{group agency; collective responsibility; corporate responsibility; ownership; discursive dilemma}

\begin{document}
\maketitle

\begin{abstract}
ACT-10-a argues that some groups are agents whose attitudes are not reducible to those of their members. We argue that such groups can also be morally responsible for what they do. An organized group can act from a mechanism that is reasons-responsive and its own: its attitudes are formed and revised through its own procedures, and it has a persisting standpoint from which to answer for its past conduct. Group responsibility does not absorb the responsibility of members, and it can be present where no member is responsible for the group's act as such.
\end{abstract}

\section{The argument}

Theorem~5.3 of ACT-10-a establishes that a group which achieves rational unity without dictatorship is an agent in the sense of its Definition~3.1, and Corollary~5.4 that its attitudes supervene on its members' only holistically. Whether such an agent is responsible is a further question. We use the account of responsible agency of ACT-01-a, as stated in Definition~2.1 of ACT-16-a: an entity is a responsible agent if it can act from a mechanism that is reasons-responsive and its own, where ownership excludes installation by others that bypasses its practical reasoning and requires a standpoint from which it takes responsibility.

\begin{proposition}\label{p:resp}
An organized group agent can be a responsible agent.
\end{proposition}

\begin{proof}
Reasons-responsiveness: a group agent forms its judgments and intentions by procedures, such as premise-based voting or deliberation in committee, that take evidence and considerations as inputs and revise outputs when the considerations change; this is the rational unity that Definition~3.1 of ACT-10-a requires, and it is a capacity to recognize and act on reasons. Ownership: the group's attitudes are installed and revised through its own decision procedures. Its members' contributions are the medium through which its practical reasoning operates, not an intervention that bypasses it, as the operations of neurons are the medium of a person's reasoning and not a manipulation of it. Standpoint: an organized group persists through changes of membership, keeps records of its decisions, can acknowledge past acts as its own, and regulates its future conduct in response to the censure of others, as when a firm revises its procedures after a public finding against it. Each condition of Definition~2.1 of ACT-16-a can therefore be met.
\end{proof}

The proof follows the conditions of \citet[ch.~7]{ListPettit2011} in substance, and extends to groups the reasoning by which \citet{French1979} took corporations to be moral persons, without his appeal to the corporate internal decision structure as a sufficient condition: a structure that merely transmitted a dictator's will would not meet the non-dictatorship principle of ACT-10-a.

\begin{proposition}\label{p:nondist}
The responsibility of a group agent for an act neither entails nor is entailed by the responsibility of its members for the same act.
\end{proposition}

\begin{proof}
By Corollary~5.4 of ACT-10-a, a group may hold attitudes that a majority, or all, of its members reject. Suppose a court of three decides a case by premise-based procedure, each judge voting sincerely and competently on each premise, and the court's verdict, which two of the three judges reject, is an injustice that the court was in a position to foresee and avoid by a different procedure. The court is responsible for the verdict, by Proposition~\ref{p:resp}. The two dissenting judges are not responsible for the verdict as such, since each voted sincerely and competently and rejected the verdict; the third is responsible for her own votes, which alone did not produce it. The court's responsibility is not the sum of theirs. Conversely, members who act wrongly within a group are responsible for their contributions whether or not the group is responsible for the outcome.
\end{proof}

Proposition~\ref{p:nondist} does not let members escape. Each member remains responsible for what she contributes, and in particular for the procedures she helps to adopt and for her failure to contest group decisions she judges wrong \citep{ListPettit2011}. Group responsibility is an addition, not a substitution.

\section{Objection and reply}

\begin{objection}
Groups act only through the bodies of their members; whatever responsibility there is attaches to those who move \citep{Velasquez1983}.
\end{objection}

\begin{reply}
That a group acts through its members shows that its agency is realized in theirs, not that it is identical to theirs. ACT-10-a shows that the group's attitudes, and so the reasons for which it acts, may differ from those of every member. Responsibility follows the reasons for which an act is done, and when these are the group's, the responsibility is the group's as well.
\end{reply}

\section{Conclusion}

Groups that are agents in the sense of ACT-10-a can satisfy the conditions of responsible agency: their mechanisms respond to reasons, are their own, and belong to a standpoint that persists. They can therefore be responsible for what they do, without diminishing the responsibility of their members, and sometimes where no member is responsible for the group's act. Such groups may also be the deployers to whom ACT-16-a assigns responsibility for the outputs of artificial systems.

\begin{thebibliography}{9}\small
\bibitem[French(1979)]{French1979} French, P.~A. (1979). The corporation as a moral person. \emph{American Philosophical Quarterly}, 16(3), 207--215.
\bibitem[List and Pettit(2011)]{ListPettit2011} List, C. and Pettit, P. (2011). \emph{Group Agency: The Possibility, Design, and Status of Corporate Agents}. Oxford University Press.
\bibitem[Velasquez(1983)]{Velasquez1983} Velasquez, M.~G. (1983). Why corporations are not morally responsible for anything they do. \emph{Business and Professional Ethics Journal}, 2(3).
\end{thebibliography}

\end{document}
